
\(\underline{M}\) a une
résolution partielle globale~:
\[
  \underline{L}_i\longrightarrow \underline{L}_{i-1}\longrightarrow\ldots\longrightarrow \underline{L}_0\longrightarrow \underline{M}\longrightarrow 0,
\]
par des \(\underline{B}\)-Modules localement libres de type fini, le noyau \(\underline{Z}_i\) de \(\underline{L}_i\to \underline{L}_{i-1}\) est encore un \(\underline{B}\)-Module de type fini (I. 2.11). Il est même de présentation finie (en faisant le même raisonnement avec \(i+1\) au lieu de \(i\)). Enfin, il est plat sur \(S\). Il vérifie donc les mêmes hypothèses que \(\underline{M}\), ce qui permet de construire de proche en proche une résolution globale graduée de \(\underline{M}\) par des \(\underline{B}\)-Modules gradués localement libres de type fini, d'après la remarque initiale. Pour \(i\geq r+1\), \(\underline{Z}_i\) est un \(\underline{B}\)-Module projectif sur tout ouvert affine (1.5.2) et par suite est localement libre en tant que \(\underline{B}\)-Module gradué, d'après 1.5.1. On peut alors, dans la résolution graduée construite, remplacer \(\underline{L}_{r+1}\) par \(\underline{Z}_r\), et les termes suivants par \(0\), d'où la résolution annoncée.

Montrons enfin le point c) de la démonstration.

\textbf{Proposition 1.7.} \uline{Soient} \(S\) \uline{un schéma quasi-compact}, \(\underline{E}\) \uline{un} \(\underline{O}_S\)\uline{-Module localement libre de type fini}, \(\underline{B}=\Sym(\underline{E})\), \uline{et} \(\underline{L}\) \uline{un} \(\underline{B}\)\uline{-Module gradué}, \uline{localement libre de type fini} (1.2). \uline{Alors il existe une filtration décroissante de} \(\underline{L}\) \uline{par des sous-}\(\underline{B}\)\uline{-Modules gradués} \(\underline{L}^{(m)}\), \uline{telle que} :
\[
\begin{array}{ll}
\text{i)}& \underline{L}^{(m)}=\underline{L}\quad\text{\uline{pour} }m\leq m_0;\\[1mm]
\text{ii)}& \underline{L}^{(m)}=0\quad\text{\uline{pour} }m\geq m_1;\\[1mm]
\text{iii)}& \text{\uline{pour tout} }m,\quad \operatorname{gr}^m(\underline{L})=\underline{L}^{(m)}/\underline{L}^{(m+1)}\\
&\quad \text{\uline{est isomorphe comme} }\underline{B}\text{\uline{-Module gradué à un}}\\
&\quad \underline{B}\text{\uline{-Module gradué de la forme} }\underline{N}_m\otimes_{\underline{O}_S}\underline{B}(m),\\
&\quad \text{\uline{où} }\underline{N}_m\text{ \uline{est un} }\underline{O}_S\text{\uline{-Module localement libre de type fini}.}
\end{array}
\]

Comme \(S\) est quasi-compact, on peut trouver des entiers \(n_0\) et \(n_1\) tels qu'au voisinage de chaque point de \(S\), les éléments d'une base homogène 
